%==============
\begin{center}
{\LARGE \bf Analog Circuit Design}
\end{center}
%==============
\bigskip

\noindent
{\bf Organizers}

\begin{itemize}

\item 
Erich Barke (barke@ims.uni-hannover.de)
\item
Ralf Sommer (rsommer@rhrk.uni-kl.de)
\end{itemize}

\medskip

\bigskip

\noindent
{\bf Description}

\medskip
Symbolic computation techniques can be effectively applied to three classes of engineering
problems in the field of analog circuit design: 

\begin{enumerate}
\item
design understanding
\item
development of generic behavioral circuit models, and 

\item
circuit sizing.
\end{enumerate}

{\em Design understanding} requires calculating circuit functions symbolically in order to determine
their functional relation to circuit behavior. 

{\em Development of generic behavioral models} is necessary for speeding up simulations of large analog
and mixed-signal circuits on system level and for numerical circuit optimization. 

{\em Circuit sizing} means the development of generic (sizable) libraries of analog building blocks. 

Supporting these tasks with computer algebra requires new algorithmic approaches which combine
symbolic computation techniques with numerical methods. Due to the extremely high complexity
of symbolic calculations, approximation techniques based on numerical magnitude information
must be used to extract dominant information from symbolic circuit functions. Hence, this session
will focus on hybrid symbolic/numeric algorithms and their applications to the design of analog
electronic circuits. 


\bigskip
\noindent
{\bf Talks}
\medskip

Date: July 19th (Friday) 

\begin{description}
\item[08:30-09:00] \ \\
  Ralf Sommer, Eckhard Hennig\\
  {\bf Overview talk on symbolic techniques for circuit analysis.} 
\item[09:30-10:00] \ \\
  Carsten Borchers, Lars Hedrich\\
  {\bf Symbolic Behavioral Model Generation and Formal Verification of Analog Circuits.
 } 
\item[09:30-10:00] \ \\
  Georges Gielen
\\
  {\bf Techniques for the symbolic analysis of distortion and intermodulation.
} 
\item[10:00-10:30] \ \\
  \\
  {\bf Coffee Break } 
\item[10:30-11:00] \ \\
  Wolfgang Mathis\\
  {\bf Nonlinear network theory, differential geometric approach to nonlinear networks.} 
\item[11:00-11:30] \ \\
  Ljiljana Trajkovic\\
  {\bf Homotopy Methods for Solving Nonlinear Equations and their Applications in Circuit
    Simulation.
} 
\item[11:30-12:00] \ \\
  B. Mourrain\\
  {\bf Integration of algebraic and numerical tools for devising efficient algorithms for the
    solution of systems of polynomial equations.} 
\end{description}

